% Figure: a heat pipe in longitudinal section, with an evaporator end that boils
% the working fluid, an adiabatic transport section where vapour flows to the cold
% end, a condenser end where the vapour condenses, and a wick that returns the
% condensate by capillary action, so the device moves heat at a near-isothermal
% temperature with an effective conductivity far above solid metal.
\begin{figure}[htbp]
\centering
\begin{tikzpicture}[>=stealth]
% outer wall
\draw[thick] (0,0) rectangle (11,2);
% wick lining top and bottom
\fill[enggray!25] (0,0) rectangle (11,0.3);
\fill[enggray!25] (0,1.7) rectangle (11,2);
\node[enggray, font=\scriptsize, anchor=west] at (4.3,0.15) {wick (capillary return)};
% vapour core
\node[engblue!70!black, font=\scriptsize] at (5.5,1.0) {vapour core};
% zone dividers
\draw[dashed, enggray] (3.3,0) -- (3.3,2);
\draw[dashed, enggray] (7.7,0) -- (7.7,2);
% zone labels
\node[font=\scriptsize, anchor=south] at (1.65,2.05) {evaporator};
\node[font=\scriptsize, anchor=south] at (5.5,2.05) {adiabatic transport};
\node[font=\scriptsize, anchor=south] at (9.35,2.05) {condenser};
% heat in at evaporator
\foreach \x in {0.6,1.2,1.8,2.4} { \draw[engred, ->, thick] (\x,-0.7) -- (\x,-0.1); }
\node[engred, font=\scriptsize, anchor=north] at (1.65,-0.75) {heat in};
% heat out at condenser
\foreach \x in {8.6,9.2,9.8,10.4} { \draw[engblue, ->, thick] (\x,2.1) -- (\x,2.7); }
\node[engblue, font=\scriptsize, anchor=south] at (9.5,2.75) {heat out};
% vapour flow arrow (evaporator to condenser, through the core)
\draw[engred!70!black, ->, very thick] (2.6,1.0) -- (8.4,1.0);
\node[engred!70!black, font=\scriptsize, anchor=south] at (5.5,1.05) {};
% liquid return arrow (condenser back to evaporator, in the wick)
\draw[engblue, ->, very thick] (8.4,0.45) -- (2.6,0.45);
\node[engblue, font=\scriptsize, anchor=north] at (5.5,0.42) {condensate return};
\end{tikzpicture}
\caption[Heat pipe in longitudinal section]{A heat pipe carries heat by
evaporating a working fluid at the hot end and condensing it at the cold end,
with a capillary wick pumping the condensate back with no moving part. Heat
entering the evaporator boils the fluid; the vapour flows down the low-resistance
core to the condenser, driven by the small pressure difference a slight
temperature drop sets up; there it gives up its latent heat and condenses; and
the wick draws the liquid back to the evaporator by capillary action to close the
loop. Because the transport is by latent heat and the vapour is nearly
isothermal, a heat pipe moves heat along its length at an effective thermal
conductivity many times that of solid copper, at a temperature drop of a few
kelvin. The binding limits are the capillary limit, where the wick can no longer
pump the condensate back against gravity and drag, and the boiling limit, where
the evaporator dries out, both of which cap the heat a given pipe can carry.}
\label{fig:vol04ch06:heat-pipe}
\end{figure}
